\section{Lösungsstrategie}

Die Anwendung wurde als Microservice-Architektur umgesetzt. Die wichtigsten
Lösungsansätze sind in Tabelle~\ref{tab:loesungsstrategie} zusammengefasst.

\begin{table}[ht]
    \centering

    {
        \renewcommand{\arraystretch}{1.5}
        \rowcolors{2}{tableRowColor}{white}

        \begin{tabularx}{\linewidth}{
                >{\bfseries\raggedright\arraybackslash}p{0.20\linewidth}
                >{\raggedright\arraybackslash}p{0.35\linewidth}
                >{\raggedright\arraybackslash}X
        }
        \toprule
        Bereich & Umsetzung & Begründung \\
        \midrule

        Systemaufteilung &
        Aufteilung in fachlich getrennte Microservices &
        Unterstützt Modularität und Erweiterbarkeit. \\

        Betrieb &
        Containerisierung mit Docker &
        Vereinfacht die Bereitstellung und lokale Ausführung der Anwendung. \\

        Zugriff &
        Zentrales API Gateway mit JWT-Validierung &
        Fasst die einzelnen Services nach außen zu einem einheitlichen System
        zusammen und stellt einen zentralen, geschützten Zugang bereit. \\

        Kommunikation &
        REST für Anfragen und MQTT für Ereignisse &
        Schafft klar definierte Kommunikationswege. \\

        Technologien &
        Spring Boot, Flask, Astro und SolidJS &
        Ermöglicht passende Technologien für jeden Bereich. \\

        \bottomrule
        \end{tabularx}
    }

    \caption{Zentrale Lösungsansätze}
    \label{tab:loesungsstrategie}
\end{table}